\documentstyle[% 


righttag% 


,11pt% 


,amssymb% 


,russian%               remove in Eng. 


,izvru%                 Set izven% in Eng. 


]{amsart} 


 


%  ... !


 


\newcommand{\Gr}{\operatorname{Gr}} 


\newcommand{\const}{\operatorname{const}} 


\newcommand{\extr}{\operatorname{extremum}} 


\newcommand{\la}{\langle} 


\newcommand{\ra}{\rangle} 


\newcommand{\mymath}[1]{\mathop{#1}\limits} 


\newcommand{\tts}[1]{} 


 


 


%\hoffset=-15mm%                  For Eng. 


\setcounter{page}{49} 


\god{2005} 


\nomer{5~(516)} %                        Remove in Eng. 


%\nomer{Vol.\,49, No.\,5}%               Set in Eng. 


%\str{\pageref{begin}--\pageref{end}}%  Set in Eng. 


\udk{514.763} 


 


\author{\it С.Г.\,Лейко} 


% NB:   ^^ remove in Eng. 


\title{Изопериметрические задачи для функционалов  поворота  


первого и второго порядков в (псевдо)римановых  


многообразиях} 


 


\date{} 


%\pagestyle{myheadings}%         Set in Eng. 


\pagestyle{plain} %              Remove in Eng. 


\begin{document} 


%\thispagestyle{plain}%          Set in Eng. 


\maketitle 


\label{begin} 


 


 


 


Работа продолжает исследование, начатое в [1], [2]. В данной статье 


доказана 


сингулярность (вырожденность) лагранжианов поворота первого и второго 


порядков в (псевдо)римановых многообразиях и рассмотрены различные 


изопериметрические задачи для этих функционалов. 


 


\section{Сингулярность функционалов поворота} 


 


Теория вариационных задач с высшими производными для 


невырожденных лагранжианов восходит к работам М.В.\,Остроградского (см. [3]). 


Попытки [4], [5] создать аналог такой теории для вариационных задач с 


вырожденными лагранжианами, которые имеют важное значение для 


приложений в теоретической физике, сталкиваются с принципиальными 


трудностями, связанными в первую очередь с тем, что уравнения 


Эйлера--Лагранжа для экстремалей не имеют нормальной формы. Вследствие 


этого 


не удается получить полного аналога уравнений Гамильтона относительно 


канонических переменных. В этом параграфе доказывается, что 


лагранжианы поворота первого и второго порядков в (псевдо)римановых 


пространствах произвольной размерности являются вырожденными. 


 


В (псевдо)римановом многообразии $(M^n,g)$ вдоль гладкой кривой 


$\gamma:\,]t_0,t_1[\to M^n$ из ее касательного вектора $\xi=\Dot\gamma$ 


построим ковариантным 


дифференцированием относительно связности Леви-Чивитa $\nabla$ 


(метрической связности без кручения ([6], с.\,153)) векторы высших 


кривизн 


$\mymath{\xi}_1=\nabla_t\xi,\dotsc,\mymath{\xi}_p=\nabla_t\mymath{\xi}_


{p-1}$. Рассмотрим функционал длины 


$\ell[\gamma]=\int\limits^{t_1}_{t_0}\sqrt{e_0\la\Dot\gamma, 


\Dot\gamma\ra}\,dt$, $e_0=\pm1$ --- знак скалярного квадрата 


(относительно метрики $g$) касательного вектора, и функционалы 


поворота $\mymath{\theta}_\alpha[\gamma]= 


\int\limits^{\ell_1}_{\ell_0}k_\alpha(\ell)d\ell$ порядка 


$\alpha=1,\dotsc,n-1$. Здесь $\mymath{k}_\alpha(\ell)$ --- кривизна 


Френе (порядка $\alpha$) данной кривой $\gamma$ ([7], c.\,480). Лагранжианы 


функционалов поворота первого и второго порядков при произвольной 


параметризации кривой имеют вид 


$$ 


L_1(x,\xi,\mymath{\xi}_1)=e_0\la\xi,\xi\ra^{-1} 


\sqrt{\varepsilon_1G_1},\quad L_2(x,\xi,\mymath{\xi}_1, 


\mymath{\xi}_2)=\varepsilon_1G^{-1}_1 


\sqrt{e_0\la\xi,\xi\ra}\,\sqrt{\varepsilon_2G_2}, 


$$ 


где $G_1=\Gr(\xi,\mymath{\xi}_1)$, $G_2=\Gr(\xi, 


\mymath{\xi}_1,\mymath{\xi}_2)$ --- определители Грама 


соответствующих систем векторов, $\varepsilon_1$, $\varepsilon_2$ --- знаки 


этих определителей. 


Применяя оператор ковариантного дифференцирования относительно 


связности Леви-Чивитa, в локальных координатах $x^i$ имеем 


$$ 


\xi^i=\Dot x{}^i,\quad \mymath{\xi}_1{}^i=\Ddot x{}^i+ 


\Gamma^i_{jk}(x)\Dot x{}^j\Dot x{}^k,\quad 


\mymath{\xi}_2{}^i=\dddot x{}^i+\partial_l\Gamma^i_{jk}(x) 


\Dot x{}^j\Dot x{}^k\Dot x{}^l+2\Gamma^i_{jk}(x) 


\Dot x{}^j\Ddot x{}^k+\Gamma^i_{jk}(x) 


\Dot x{}^j\mymath{\xi}_1{}^k,\dotsc 


$$ 


Отсюда видно, что выражения для лагранжианов поворота 


$L_\alpha(x,\xi,\mymath{\xi}_1,\dotsc,\mymath{\xi}_\alpha)$ любого порядка 


$\alpha$ включают высшие старшие 


производные порядка $\alpha+1$ локальных координат $x^i$ многообразия. 


Кроме 


того, $\frac{\partial L_\alpha}{\partial\mymath{x^i}^{(\alpha)}}= 


\frac{\partial L_\alpha}{\partial\mymath{\xi}_\alpha{}^i}$. 


 


\begin{thm} 


Лагранжианы поворота $L_\alpha(x,\xi,\mymath{\xi}_1,\dotsc, 


\mymath{\xi}_\alpha)$ первого и 


второго порядков являются сингулярными $($вырожденными$)$, т.\,е. для них 


старший гессиан $\det\big(\frac{\partial^2L_\alpha} 


{\partial\mymath{\xi}_\alpha{}^i\partial\mymath{\xi}_\alpha{}^j} 


\big)\equiv0$. 


\end{thm} 


 


\begin{pf} 


Достаточно показать, что матрицы гессианов 


имеют собственные векторы, отвечающие нулевому собственному 


значению. 


 


1) Рассмотрим сначала лагранжиан поворота первого порядка. Для 


определителя Грама 


$$ 


G_1(x,\xi,\mymath{\xi}_1)= 


\begin{vmatrix} 


\la\xi,\xi\ra & \la\xi,\xi_1\ra \\ 


\la\xi_1,\xi\ra & \la\xi_1,\xi_1\ra 


\end{vmatrix} 


=g_{ij}\xi^i\xi^j g_{kl}\mymath{\xi}_1{}^k 


\mymath{\xi}_1{}^l-(g_{ij}\xi^i\xi^j)^2 


$$ 


находим частным дифференцированием 


$$ 


\frac{\partial G_1}{\partial\mymath{\xi}_1{}^i}= 


-2\la\xi,\xi_1\ra\xi_i+2\la\xi,\xi\ra\mymath{\xi}_1{}_i, 


\ \ \xi_i=g_{ij}\xi^j,\ \ \mymath{\xi}_1{}_i=g_{ij} 


\mymath{\xi}_1{}^j,\ \ 


\frac{\partial^2G_1}{\partial\mymath{\xi}_1{}^j 


\partial\mymath{\xi}_1{}^i}=-2\xi_i\xi_j+2\la\xi,\xi\ra g_{ij}. 


$$ 


Отсюда имеем свертки 


$$ 


\frac{\partial G_1}{\partial\mymath{\xi}_1{}^i}\xi^i=0, 


\ \ \frac{\partial^2G_1}{\partial\mymath{\xi}_1{}^j 


\partial\mymath{\xi}_1{}^i}\xi^j=0,\ \ 


\frac{\partial G_1}{\partial\mymath{\xi}_1{}^i} 


\mymath{\xi}_1{}^i=2G_1,\ \ 


\frac{\partial^2G_1}{\partial\mymath{\xi}_1{}^i 


\partial\mymath{\xi}_1{}^j}\mymath{\xi}_1{}^j= 


\frac{\partial G_1}{\partial\mymath{\xi}_1{}^i}. 


$$ 


Находим элементы матрицы гессиана 


$$ 


\frac{\partial^2L_1}{\partial\mymath{\xi}_1{}^j 


\partial\mymath{\xi}_1{}^i}=\frac{e_0\varepsilon_1} 


{2\la\xi,\xi\ra}\bigg[-\frac{1}{2}\varepsilon_1(\varepsilon_1G_1 


)^{-3/2}\frac{\partial G_1}{\partial\mymath{\xi}_1{}^j} 


\frac{\partial G_1}{\partial\mymath{\xi}_1{}^i}+ 


(\varepsilon_1G_1)^{-1/2} 


\frac{\partial^2G_1}{\partial\mymath{\xi}_1{}^j 


\partial\mymath{\xi}_1{}^i}\bigg] 


$$ 


и, используя полученные свертки, получаем 


$$ 


\frac{\partial^2L_1}{\partial\mymath{\xi}_1{}^j 


\partial\mymath{\xi}_1{}^i}\xi^j=0,\quad 


\frac{\partial^2L_1}{\partial\mymath{\xi}_1{}^j 


\partial\mymath{\xi}_1{}^i}\mymath{\xi}_1{}^j=0. 


$$ 


 


Таким образом, имеем два собственных вектора $\xi^i$, $\mymath{\xi}_1{}^i$ 


матрицы старшего гессиана, отвечающих нулевому собственному 


значению. В частности, векторы $\xi^i,\eta^i=\la\xi,\xi_1\ra\xi^i- 


\la\xi,\xi\ra\mymath{\xi}_1{}^i$ 


образуют ортогональную пару таких собственных векторов. 


 


2) Рассмотрим теперь лагранжиан поворота второго порядка. Для него 


\begin{gather*} 


\frac{\partial^2L_2}{\partial\mymath{\xi}_2{}^j 


\partial\mymath{\xi}_2{}^i}=\frac{\varepsilon_1\varepsilon_2}{2G_1} 


\sqrt{e_0\la\xi,\xi\ra}\bigg[-\frac{\varepsilon_2}{2} 


(\varepsilon_2G_2)^{-3/2}\frac{\partial G_2} 


{\partial\mymath{\xi}_2{}^j}\frac{\partial G_2} 


{\partial\mymath{\xi}_2{}^i}+(\varepsilon_2G_2)^{-1/2} 


\frac{\partial^2G_2}{\partial\mymath{\xi}_2{}^j 


\partial\mymath{\xi}_2{}^i}\bigg],\\ 


% 


G_2(x,\xi,\mymath{\xi}_1,\mymath{\xi}_2)= 


\begin{vmatrix} 


\la\xi,\xi\ra & \la\xi,\xi_1\ra & \la\xi,\xi_2\ra \\ 


\la\xi_1,\xi\ra & \la\xi_1,\xi_1\ra & \la\xi_1,\xi_2\ra \\ 


\la\xi_2,\xi\ra & \la\xi_2,\xi_1\ra & \la\xi_2,\xi_2\ra 


\end{vmatrix} 


,\\ 


\frac{\partial G_2}{\partial\mymath{\xi}_2{}^i}=\xi^i( 


\la\xi,\xi_1\ra \la\xi_1,\xi_2\ra-\la\xi,\xi_2\ra 


\la\xi_1,\xi_1\ra)+\la\xi,\xi_2\ra(\la\xi,\xi_1\ra 


\mymath{\xi}_1{}_i-\la\xi_1,\xi_1\ra\xi_i)-\\ 


% 


-\mymath{\xi}_1{}_i(\la\xi,\xi\ra \la\xi_1,\xi_2\ra- 


\la\xi,\xi_1\ra \la\xi,\xi_2\ra)-\la\xi_1,\xi_2\ra 


(\la\xi,\xi\ra\mymath{\xi}_1{}_i-\la\xi,\xi_1\ra\xi_i+ 


2G_1\mymath{\xi}_2{}_i,\\ 


% 


\frac{\partial^2G_2}{\partial\mymath{\xi}_2{}^j 


\partial\mymath{\xi}_2{}^i}=\xi_i(\la\xi,\xi_1\ra 


\mymath{\xi}_1{}_j-\la\xi_1,\xi_1\ra\xi_j)+\xi_j(\la\xi,\xi_1\ra 


\mymath{\xi}_1{}_i-\la\xi_1,\xi_1\ra\xi_i)- \\ 


% 


-\mymath{\xi}_1{}_i(\la\xi,\xi\ra\mymath{\xi}_1{}_j- 


\la\xi,\xi_1\ra\xi_j)-\mymath{\xi}_1{}_j(\la\xi,\xi\ra 


\mymath{\xi}_1{}_i-\la\xi,\xi_1\ra\xi_i)+2G_1g_{ij}. 


\end{gather*} 


Свертки дают следующие значения: 


\begin{gather*} 


\frac{\partial G_2}{\partial\mymath{\xi}_2{}^i}=0,\quad 


\frac{\partial^2G_2}{\partial\mymath{\xi}_2{}^j 


\partial\mymath{\xi}_2{}^i}\xi^j=0,\quad 


\frac{\partial G_2}{\partial\mymath{\xi}_2{}^i} 


\mymath{\xi}_1{}^i=0,\\ 


% 


\frac{\partial^2G_2}{\partial\mymath{\xi}_2{}^j 


\partial\mymath{\xi}_2{}^i}\mymath{\xi}_1{}^i=0,\quad 


\frac{\partial G_2}{\partial\mymath{\xi}_2{}^i} 


\mymath{\xi}_2{}^i=2G_2,\quad 


\frac{\partial^2G_2}{\partial\mymath{\xi}_2{}^j 


\partial\mymath{\xi}_2{}^i}\mymath{\xi}_2{}^j= 


\frac{\partial G_2}{\partial\mymath{\xi}_2{}^i}. 


\end{gather*} 


Используя эти значения, убеждаемся в том, что векторы $\xi^i$, 


$\mymath{\xi}_1{}^i$, $\mymath{\xi}_2{}^i$ 


являются собственными векторами матрицы гессиана $L_2$, 


отвечающими нулевому собственному значению. В частности, векторы 


\begin{multline*} 


\xi^i,\eta^i,\zeta^i=(\la\xi,\xi_1\ra \la\xi_1,\xi_2\ra- 


\la\xi,\xi_2\ra \la\xi_1,\xi_1\ra)\mymath{\xi}_2{}^i+ \\ 


% 


+(\la\xi,\xi\ra \la\xi_1,\xi_2\ra- \la\xi,\xi_1\ra 


\la\xi,\xi_2\ra)\mymath{\xi}_1{}^i+(\la\xi,\xi\ra 


\la\xi_1,\xi_1\ra-\la\xi,\xi_2\ra^2)\xi^i 


\end{multline*} 


образуют ортогональную тройку таких собственных векторов. 


\end{pf} 


 


 


По-видимому, обнаруженная ситуация будет сохраняться и в 


случае лагранжиана поворота произвольного порядка $\alpha=1,\dotsc,n-1$, 


для которого касательный вектор и векторы кривизн до порядка $\alpha$


включительно являются собственными векторами, отвечающими 


нулевому собственному значению. 


 


 


\section{Изопериметрические задачи} 


 


Исследование свойств экстремалей функционалов поворота с точки 


зрения дифференциальной геометрии (псевдо)римановых 


многообразий представляет интерес в рамках общей теории 


вырожденных лагранжианов, а также для приложений. Так, например, 


в [4] рассматривался случай, когда обобщенное действие для точечной 


релятивистской частицы включает функционал длины и функционал 


поворота первого порядка 4-мерного пространства Минковского. При 


этом собственные векторы матрицы гессиана, отвечающие нулевому 


собственному значению, определяют первичные лагранжевы связи 


теории таких частиц. Обозначим через 


\begin{gather*} 


\delta\ell_i=\frac{\partial L}{\partial x^i}-\frac{d}{dt} 


\frac{\partial L}{\partial\Dot x{}^i},\quad 


\delta\mymath{\theta}_1{}_i=\frac{\partial L_1}{\partial x^i}- 


\frac{d}{dt}\frac{\partial L_1}{\partial\Dot x{}^i}+ 


\frac{d^2}{dt^2}\frac{\partial L_1}{\partial\Ddot x{}^i},\\ 


\delta\mymath{\theta}_2{}_i=\frac{\partial L_2}{\partial x^i}- 


\frac{d}{dt}\frac{\partial L_2}{\partial\Dot x{}^i}+ 


\frac{d^2}{dt^2}\frac{\partial L_2}{\partial\Ddot x{}^i}- 


\frac{d^3}{dt^3}\frac{\partial L_2}{\partial\dddot{x}^i}  %%%% ... ?


\end{gather*} 


вариационные производные функционалов длины и поворота первого 


и второго порядков. 


 


Поскольку $\delta\ell_i=-e_0\nabla_t(e_0\la\xi,\xi\ra)^{-1/2} 


\xi_i-e_0(e_0\la\xi,\xi\ra)^{-1/2}\mymath{\xi}_1{}_i$ 


(здесь справа ковариантные компоненты векторов), отсюда имеем 


$\delta\ell_i\xi^i=0$. Таким образом, функционал длины также сингулярный. 





Уравнения Эйлера--Лагранжа $\delta\ell_i=0$, 


$\delta\mymath{\theta}_\alpha{}_i=0$ определяют 


соответственно геодезические кривые и свободные (безусловные) 


экстремали поворота (СЭП). В случае, когда варьированные кривые в 


двухточечной вариационной задаче подчинены дополнительному 


интегральному условию, стационарные кривые функционала длины и 


поворота называем изопериметрическими экстремалями длины (ИЭД) 


и, соответственно --- изопериметрическими экстремалями поворота 


(ИЭП). При варьировании надо не допускать изотропных направлений 


$\la\xi,\xi\ra=0$ и точек уплощения $(n-1)$-го порядка ($G_1=0$) для 


поворота первого порядка, и точек уплощения $(n-1)$-го и $(n-2)$-го 


порядков ($G_1=G_2=0$) для поворота второго порядка. 


 


Геодезические кривые имеют максимальное уплощение порядка 


$n-1$. Они доставляют функционалу поворота первого порядка 


минимальное нулевое значение. Их естественно называть 


тривиальными СЭП или ИЭП первого порядка. 


 


Кривые с уплощением ([7]; [8], с.\,125) порядка $n-2$ (2-геодезические 


кривые) имеют параллельную соприкасающуюся плоскость $E_2(\xi,\xi_1)$. 


Вдоль этих кривых выполняются дифференциальные уравнения вида 


$\mymath{\xi}_2{}^i=a(t)\xi^i+a_1(t)\mymath{\xi}_1{}^i$ 


(справа некоторые функции параметра на кривой). Они доставляют 


функционалу поворота второго порядка минимальное нулевое 


значение. Их естественно называть тривиальными СЭП или ИЭП 


второго порядка. 


 


\begin{thm} 


Для первого поворота пучки ИЭП и ИЭД совпадают. 


\end{thm} 


 


\begin{pf} 


В работе [1] были исследованы СЭП 


и ИЭП первого порядка при условии фиксирования длины, т.\,е. 


экстремали вариационной задачи 


$$ 


\extr\mymath{\theta}_1[\gamma], \quad\ell[\gamma]=\mymath{\ell}^\frown 


\text{ --- } \const. 


$$ 


Их дифференциальные уравнения имеют вид $c_0\delta 


s_i+c_1\delta\mymath{\theta}_1{}_i=0$, $c_0,\ c_1$ --- $\const$. 


В случаях $c_0=0$ или $c_1=0$ получаем уравнения СЭП первого порядка 


или геодезическую кривую (тривиальную ИЭП). В остальных случаях 


уравнения ИЭП первого порядка включают одну существенную 


постоянную, которая зависит от фиксированной длины. 


 


Оказывается, что указанные выше дифференциальные уравнения 


имеют также ИЭД при условии фиксирования первого поворота 


первого порядка среди всех дуг с фиксированными концами $p_0$, $p_1$. 


Действительно, рассмотрим двухточечную вариационную задачу 


$$ 


\extr\ell[\gamma],\quad \mymath{\theta}_1[\gamma]= 


\mymath{\theta}^\frown_1\text{ --- } \const. 


$$ 


Допустим, что $x^h=\mymath{x}^\frown{}^h(t)$ --- параметрические уравнения 


кривой $\overset{\frown}{\gamma}$, 


на которой функционал длины достигает экстремального значения. 


 


Проварьируем эту кривую пучком допустимых кривых 


$$ 


\gamma:x^h(t,\varepsilon_1,\varepsilon_2)= 


\mymath{x}^\frown{}^h(t)+\varepsilon_1\eta^h_1(t)+ 


\varepsilon_2\eta^h_2(t). 


$$ 


Здесь $\varepsilon_1$, $\varepsilon_2$ --- малые параметры, $\eta^h_1(t)$, 


$\eta^h_2(t)$ --- гладкие вектор-функции, 


которые удовлетворяют на концах касанию первого порядка: 


$$ 


\eta^h_{1,2}(t_0)=\eta^h_{1,2}(t_1)=\Dot\eta{}^h_{1,2}(t_0)= 


\Dot\eta{}^h_{1,2}(t_1)=0. 


$$ 


На этот пучок в данной задаче наложено интегральное условие 


постоянства первого поворота 


$$ 


F(\varepsilon_1,\varepsilon_2)=\mymath{\theta}_1[\gamma]= 


\mymath{\theta}_1[\overset{\frown}{\gamma}{}+\varepsilon_1\eta_1+ 


\varepsilon_2\eta_2]=\mymath{\theta}_1^\frown. 


$$ 


Из него находим 


$$ 


F_{\varepsilon_2}(0,0)=\frac{\partial F(0,0)} 


{\partial\varepsilon_2}=\delta\mymath{\theta}_1[\eta_2]= 


\int^{t_1}_{t_0}\delta\mymath{\theta}_1{}_i\eta^i_2dt. 


$$ 


 


В случае $\delta\mymath{\theta}_1[\eta_2]=0$ для всех указанных допустимых 


вектор-функций $\eta^h_2(t)$ 


получаем $\delta\mymath{\theta}_1{}_i=0$. 


 


Допустим, что существует вектор-функция $\eta^h_2(t)$ такая, что 


$\delta\mymath{\theta}_1[\eta_2]\ne0$. В этом случае 


$F_{\varepsilon_2}(0,0)\ne0$. Тогда к уравнению 


$F(\varepsilon_1,\varepsilon_2)=\mymath{\theta}^\frown_1$ применима теорема 


о неявной функции. На основании 


этой теоремы данное уравнение локально разрешимо в классе гладких 


функций, т.\,е. существует функция $\varepsilon_2(\varepsilon_1)$, 


$\varepsilon_2(0)=0$ такая, что 


$F(\varepsilon_1,\varepsilon_2(\varepsilon_1))= 


\mymath{\theta}^\frown_1$ для всех достаточно малых значений 


$\varepsilon_1$. Кроме того, 


$$ 


\varepsilon'_2(0)=-\frac{F_{\varepsilon_1}(0,0)} 


{F_{\varepsilon_2}(0,0)}=-\frac{\delta\mymath{\theta}_1[\eta_1]} 


{\delta\mymath{\theta}_1[\eta_2]}. 


$$ 


Поскольку по постановке задачи на кривой $\overset{\frown}{\gamma}+ 


\varepsilon_1\eta_1+\varepsilon_2(\varepsilon_1)\eta_2$ 


функционал $\ell[\gamma]$ должен достигать экстремального значения, то 


необходимо 


$$ 


\frac{d}{d\varepsilon_1}\ell[\mymath{\gamma}^\frown+ 


\varepsilon_1\eta_1+\varepsilon_2(\varepsilon_1)\eta_2 


]\Big|_{\varepsilon_2=0}=\delta\ell[\eta_1]+ 


\varepsilon'_2(0)\delta\ell[\eta_2]=0. 


$$ 


 


Введем обозначение $\mymath{c}^\frown=-\delta\ell[\eta_2]/ 


\delta\mymath{\theta}_1[\eta_2]$ 


и запишем предыдущее равенство в виде $\delta\ell[\eta_1]+ 


\mymath{c}^\frown \delta\mymath{\theta}_1[\eta_1]=0$. 


Последнее равенство имеет тождественный характер для всех гладких 


вектор-функций $\eta^h_1(t)$, удовлетворяющих на концах касанию первого 


порядка. Поэтому $\delta\ell_i+\mymath{c}^\frown\delta 


\mymath{\theta}_1{}_i=0$. 


 


Таким образом, пучки ИЭП и ИЭД, для которых функционал первого 


поворота и изопериметрические условия являются взаимными, 


совпадают. 


\end{pf} 


 


 


Для поворота второго порядка можно ставить как безусловную 


вариационную задачу, так и различные изопериметрические задачи при 


условии постоянства длины и (или) первого поворота. Применяя 


рассуждения предыдущей теоремы (или леммы об изометрических 


постоянных [1], [9]), можно получить соответствующие 


дифференциальные уравнения экстремалей. Например, для 


функционала поворота второго порядка изопериметрическая задача 


$$ 


\extr\mymath{\theta}_2[\gamma],\quad 


\mymath{\theta}_1[\gamma]=\mymath{\theta}_1^\frown 


\text{ --- }\const,\quad \ell[\gamma]= 


\mymath{\ell}^\frown\text{ --- }\const 


$$ 


порождает пучок ИЭП второго порядка с дифференциальными 


уравнениями вида $c_2\delta\mymath{\theta}_2{}_i+ 


c_1\delta\mymath{\theta}_1{}_i+ +c\delta\ell_i=0$,


%% remove extra + in Eng.     ^^^


$c_1,\ c_2,\ c$ --- $\const$. 


При $c=c_1=0$ получим уравнения СЭП второго порядка 


$\delta\mymath{\theta}_2{}_i=0$. 


Как и в случае теоремы 2, можно показать, что последние уравнения 


описывают пучки изопериметрических экстремалей для длины или 


первого поворота с взаимными изопериметрическими условиями на 


первый и второй повороты или длину и второй поворот. 


 


Вариационные производные функционала первого (второго) 


поворота включают векторы кривизны третьего (пятого) порядка. 


Записывая явные выражения этих производных (их можно найти в 


[1], [10]), получим дифференциальные уравнения 


соответствующих экстремалей. Эти уравнения оказываются 


инвариантными относительно выбора параметра кривой и выбора 


локальных координат многообразия. Вид этих уравнений имеет 


достаточно сложный характер и, что следует особо подчеркнуть, они не 


имеют в общем случае нормальной формы. 


 


 


\section{Изопериметрические задачи для функционала поворота на 


поверхностях} 


 


В случае двумерного риманова многообразия 


$(M^2,g)$ определен только один поворот (первый). Поскольку кривые 


на многообразиях минимальной размерности имеют уплощение не 


выше первого порядка, то дифференциальные уравнения 


изопериметрических экстремалей поворота значительно упрощаются. 


Как было показано [1], эти уравнения имеют яркий геометрический 


характер --- вдоль ИЭП геодезическая и гауссовы кривизны 


пропорциональны: $k_1=cK$, где $c$ --- некоторая постоянная (зависящая 


от фиксированной длины экстремали) и $K$ --- гауссова кривизна 


многообразия. Примечательно, что здесь можно получить уравнения 


ИЭП и в нормальной форме. Именно, на областях многообразия, где 


гауссова кривизна ненулевая и сохраняет знак, последнее уравнение 


эквивалентно системе дифференциальных уравнений вида 


$\nabla_\ell\mymath{\xi}_1{}^i=\la\xi,\xi\ra^{-1} 


\la\xi_1,\xi_1\ra\xi^i+(K^{-1}\nabla_\ell K) \mymath{\xi}_1{}^i$ 


(здесь дифференцирование ведется по длине $\ell$ дуги кривой) с 


начальными условиями $\la\xi(0),\xi(0)\ra=1$, $\la\xi(0),\xi_1(0)\ra=0$. 


 


Эти уравнения имеют нормальную форму и применение к ним 


стандартных теорем из теории дифференциальных уравнений позволяет 


решать для ИЭП задачу Коши [11]. 


 


К уравнениям изопериметрических экстремалей поворота в 


форме $k_1=cK$ приводят и другие вариационные задачи. Так, при 


изучении А.\,Пуанкаре астрономической задачи ``о трех телах'' возникла 


необходимость отыскания замкнутых геодезических кривых на 


овальной поверхности. С этой целью рассматривалась 


изопериметрическая задача 


$$ 


\min\ell[\gamma],\quad \Omega[G]=\iint\limits_G 


K\sqrt{\det g}\,dx^1\wedge dx^2=\mymath{\Omega}^\frown=\const, 


$$ 


где $\gamma$ --- замкнутая кривая, ограничивающая область $G$ поверхности, 


$\Omega$ --- 


интегральная гауссова кривизна этой области. Установлено ([12], с.\,229), 


что стационарные кривые этой задачи описываются уравнениями 


$k_1=cK$. 


 


Задача об отыскании замкнутой линии поверхности, 


охватывающей наибольшую площадь при заданной длине, 


рассматривалась Ф.\,Миндингом ([13], с.\,128). Экстремалями этой задачи 


служат кривые постоянной геодезической кривизны (окружности 


Дарбу). Таким образом, на поверхностях постоянной гауссовой 


кривизны окружности Дарбу являются изопериметрическими 


экстремалями поворота. 


 


ИЭП естественным образом возникают в дифференциальной 


геометрии касательных расслоений. Для того чтобы кривая в 


сферическом касательном расслоении $T_\rho(M^2)$ с метрикой Сасаки 


$g^S$ двумерного неплоского риманова многообразия $(M^2,g)$ была 


геодезической, необходимо и достаточно, чтобы ее проекция на 


базисное многообразие $M^2$ была изопериметрической экстремалью 


поворота [14], [15]. 


 


На поверхностях с нулевой гауссовой кривизной (торсах) 


двухточечная безусловная, а также изопериметрическая задача для 


функционала поворота имеют мощные пучки решений. На таких 


поверхностях всякая гладкая кривая является стационарной кривой 


вариационной задачи. 


 


Вариационные задачи для функционала поворота на поверхностях 


можно ставить посредством использования теоремы Гаусса--Боннэ. 


Пусть в евклидовой плоскости $E^2$ с декартовыми прямоугольными 


координатами $x$, $y$ задана гладкая без самопересечений кривая $\gamma_0$, 


ограничивающая область $D\subset E^2$. Будем искать вложение 


$f:D\to E^3$ такое, чтобы на вложенной поверхности $f(D)$ кривая 


$\gamma=f(\gamma_0)$ имела экстремальный поворот. Посредством формулы 


Гаусса--Боннэ $\theta[\gamma]+\Omega[D]=2\pi$ 


эта задача приводит к отысканию экстремальных поверхностей 


функционала интегральной гауссовой кривизны 


$$ 


\Omega[f]=\iint\limits_D K\sqrt{\det g}\,dx\wedge dy. 


$$ 


Здесь $K=(\det g)^{-1/2}(f_{xx}f_{yy}-f_{xy}^2)$, $\det g=1+f^2_x+f^2_y$. 


 


 


\begin{thm} 


Всякая гладкая поверхность $z=f(x,y)$ класса 


$C^4(D)$ является стационарной поверхностью для функционала 


интегральной полной кривизны. 


\end{thm} 


 


\begin{pf} 


Лагранжиан $L=K\sqrt{f}$ функционала $\Omega[f]$ 


зависит от вторых производных графика $f$ поверхности $z=f(x,y)$. 


Как известно ([9], с.\,205), стационарные поверхности в этом случае 


удовлетворяют уравнениям Эйлера--Пуассона 


$$ 


\frac{\delta\Omega[f]}{\delta f}= 


\frac{\partial L}{\partial f}-\frac{\partial}{\partial x} 


\frac{\partial L}{\partial f_x}-\frac{\partial}{\partial y} 


\frac{\partial L}{\partial f_y}+ 


\frac{\partial^2}{\partial x\partial y} 


\frac{\partial L}{\partial f_{xy}}+\frac{\partial^2}{\partial x^2} 


\frac{\partial L}{\partial f_{xx}}+ 


\frac{\partial^2}{\partial y^2} 


\frac{\partial L}{\partial f_{yy}}=0. 


$$ 


Непосредственная проверка показывает (мы опускаем из-за 


громоздкости), что эти уравнения удовлетворяются для всякой 


функции $f\in C^4(D)$. 


\end{pf} 


 


Мощность пучка стационарных кривых функционала 


интегральной кривизны имеет яркое геометрическое толкование. 


Действительно, поверхности, касающиеся на кривой $\gamma=f(\gamma_0)$, 


доставляют ей одинаковый поворот, поскольку из геометрического 


смысла геодезической кривизны кривой на поверхности вытекает, что 


геодезическая кривизна кривой $\gamma=f(\gamma_0)$ на всех касающихся 


варьируемых допустимых поверхностях одинакова. Здесь 


прослеживается интересная аналогия с тем, что стационарные кривые 


для функционала поворота на торсах, имеющие на концах касание 


первого порядка, имеют с точностью до кратных $2\pi$ одинаковый 


поворот. Кроме того, теорема 3 показывает, что для функционала 


интегральной кривизны поверхности вариационная задача отыскания 


стационарных поверхностей с экстремальным поворотом края требует 


дополнительных изопериметрических или других ограничений 


геометрического содержания. Соответствующие задачи могут быть 


предметом отдельного исследования. 


 


Отметим также полученную нами механическую интерпретацию 


для изопериметрических экстремалей поворота, которые описывают 


равновесные состояния однородной, нерастяжимой нити в отсутствие 


поверхностного трения на поверхности ненулевой гауссовой кривизны 


[16]. 


 


\begin{thebibliography}{99} 


 


\bibitem{1} 


Лейко С.Г. {\em Вариационные задачи для функционалов поворота и 


спин-отображения псевдоримановых пространств} // Изв. вузов. Математика. -- 


1990. -- \No\,10. -- С.\,9--17. 


\bibitem{2} 


Лейко С.Г. {\em Экстремали функционалов поворота кривых 


псевдориманова пространства и траектории спин-частиц в 


гравитационных полях} // Докл. РАН. -- 1992. -- Т.325. 


-- \No\,4. -- С.\,659--664. 


\bibitem{3} 


Остроградский М.В. {\em Мемуар о дифференциальных уравнениях, 


относящихся к изопериметрической задаче} // Полное собрание трудов. -- 


Киев: Изд-во АН УССР, 1961. -- Т.\,2. -- С.\,139--233.


\bibitem{4} 


Нестеренко В.В. {\em Формализм Остроградского для сингулярных 


лагранжианов с высшими производными}. -- В кн. Mykhailo 


Ostrohrads'kyi. Honoring his bicentenary. Institute of Mathematics. Kyiv, 


2001. -- P.\,82--99. 


\bibitem{5} 


Гриффитс Ф. {\em Внешние дифференциальные системы и вариационное 


исчисление}. -- М.: Мир, 1986. -- 360~с. 


\bibitem{6} 


Норден А.П. {\em Пространства аффинной связности}. -- М.: Наука, 1976. 


-- 432 с. 


\bibitem{7} 


Рашевский П.К. {\em Риманова геометрия и тензорный анализ}. -- М.: Наука, 


1967. -- 664 с. 


\bibitem{8} 


Лейко С.Г. {\em Рiманова геометрiя. Навчальний посiбник}. -- Одеса.: 


Астропринт, 1999. -- 211 с. 


\bibitem{9} 


Коша А. {\em Вариационное исчисление}. -- М.: Высш. школа, 1983. -- 279~с. 


\bibitem{10} 


Лейко С.Г., Мосяндз В.А. {\em Экстремали поворота $2$-го порядка в 


римановом пространстве и траектории спин-частиц в гравитационном 


поле} // Гравитация и теория относительности. Казань: Изд-во Казанск.


ун-та, 1992. -- Вып.\,29. -- С.\,72--80. 


\bibitem{11} 


Лейко С.Г. {\em Изопериметрические экстремали поворота на 


поверхностях в евклидовом пространстве $E^3$} // Изв. вузов. 


Математика. -- 1996. -- \No\,6. -- С.\,25--32. 


\bibitem{12} 


Бляшке В.В. {\em Введение в дифференциальную геометрию}. -- М.: Наука, 


1973. -- 440~с. 


\bibitem{13} 


Бляшке В.В. {\em Дифференциальная геометрия}. -- М.--Л.: ОНТИ--НКТП 


СССР, 1935. -- 332~с. 


\bibitem{14} 


Nagy P. {\em On the tangent sphere bundle of Riemannian $2$-manifolds} // 


Tohoku Mat. J. -- 1977. -- V.\,29. -- P.\,203--208. 


\bibitem{15} 


Лейко С.Г. {\em О геодезическом потоке на сферическом касательном 


расслоении двумерного многообразия с метрикой Сасаки} // 


Изв. вузов. Математика. -- 2001. -- \No\,3. -- С.\,33--38. 


\bibitem{16} 


Лейко С.Г. {\em Механическая интерпретация изопериметрических 


экстремалей поворота на поверхностях} // Вiсник Одеськ. держ. 


ун-ту. -- 1999. -- Т.\,4. -- \No\,4. -- С.\,79--82. 


 


\end{thebibliography} 


 


\vskip 2cm 


 


\post{\it Одесский национальный}{\it Поступила} 


\post{\it университет}{10.06.2004} 


 


\label{end} 


\end{document} 


